\documentclass[11pt]{article}
\usepackage[a4paper,margin=1in]{geometry}
\usepackage{amsmath,amssymb,amsthm}
\usepackage[T1]{fontenc}
\usepackage{lmodern}
\usepackage[hidelinks]{hyperref}
\setlength{\emergencystretch}{3em}

\theoremstyle{plain}
\newtheorem{theorem}{Theorem}
\newtheorem{lemma}[theorem]{Lemma}
\theoremstyle{definition}
\newtheorem{definition}[theorem]{Definition}
\theoremstyle{remark}
\newtheorem{remark}[theorem]{Remark}

\title{A Counterexample to Conjecture 00000002058:\\
Matrix Algebras Are Equationally Noetherian}
\author{%
  Feiyu\thanks{Independent researcher.}
  \and
  Claude (Opus 5.5)\thanks{Anthropic.}%
}
\date{2026-10-02}

\begin{document}
\maketitle

\begin{abstract}
Conjecture 00000002058 asserts that algebraic sets over the matrix algebra
$M_n(K)$ are Noetherian if and only if $n\le2$, and in particular that $M_3(K)$
is not equationally Noetherian. In fact $M_n(K)$ is equationally Noetherian for
every $n$ and every field $K$. Matrix equations are, entrywise, polynomial
equations in the entries of the unknowns, so Hilbert's basis theorem applies.
The argument is formalized in Lean~4 against Mathlib.
\end{abstract}

\section{The conjecture as stated}

The original text of conjecture 00000002058 reads:

\begin{quote}
Definition: Equational Noetherianity means every finite system of equations is
equivalent to a finite subsystem. Conjecture: Algebraic sets over $M_n(K)$ are
Noetherian if and only if $n \le 2$; concretely, algebraic sets over $M_3(K)$ are
non-Noetherian (by a construction of infinite ascending chains of commutator
polynomials).
\end{quote}

\section{Reading of the statement}

The definition as printed (``every \emph{finite} system \ldots{} is equivalent
to a finite subsystem'') is satisfied by every algebra, since a finite system
is a finite subsystem of itself. Read that way, $M_3(K)$ is trivially
Noetherian and the conjecture is false. We therefore use the standard notion
from algebraic geometry over algebraic structures \cite{bmr}, which is
stronger and makes the statement meaningful.

\begin{definition}
An \emph{equation} over $M_n(K)$ in unknowns $X_1,\ldots,X_k$ is a
noncommutative polynomial in the $X_l$ with coefficients in $K$ and constants
from $M_n(K)$; it is evaluated at a tuple of matrices in the obvious way. A system
$S$ of equations (of any cardinality) defines the \emph{algebraic set}
$V(S)\subseteq M_n(K)^k$ of its common zeros. $M_n(K)$ is \emph{equationally
Noetherian} if for every $k$ and every system $S$ there is a finite $S_0\subseteq
S$ with $V(S_0)=V(S)$; equivalently, algebraic sets satisfy the descending chain
condition.
\end{definition}

The conjecture asserts: for every field $K$, $M_n(K)$ is equationally
Noetherian if and only if $n\le2$.

\section{Result}

\begin{theorem}\label{thm:main}
For every field $K$ and every $n$, $M_n(K)$ is equationally Noetherian. In
particular $M_3(\mathbb Q)$ is, and Conjecture 00000002058 is false.
\end{theorem}

\section{Proof}

Let $R=K[x_{l,i,j}]$ be the polynomial ring in $kn^2$ variables, one for each
entry of each unknown. Substituting the \emph{generic matrices}
$(x_{l,i,j})_{i,j}$ for the unknowns and $(c_{ij})\in M_n(R)$ for each constant
$c\in M_n(K)$ gives an algebra homomorphism $G$ from equations to $M_n(R)$. For
every tuple of matrices $X$, the value of an equation $p$ at $X$ is obtained by
evaluating the entries of $G(p)$ at the entries of $X$: both sides are algebra
homomorphisms of $p$ that agree on unknowns and constants. Hence
\[
  V(S)=\{X : G(p)_{ij}(X)=0\ \text{for all } p\in S,\ i,j\}.
\]
$R$ is Noetherian (Hilbert's basis theorem), so the ideal generated by the
polynomials $G(p)_{ij}$ ($p\in S$) is generated by finitely many of them, coming
from finitely many equations $p_1,\ldots,p_r\in S$. If $X$ is a common zero of
$p_1,\ldots,p_r$, then evaluation at $X$ kills these generators, hence the whole
ideal, hence every $G(p)_{ij}$. So $V(\{p_1,\ldots,p_r\})=V(S)$.

\section{Formalization}

The accompanying Lean~4 project (\texttt{lean/Submission/Basic.lean}) states
the conjecture and proves its negation against Mathlib.

\begin{itemize}
  \item \texttt{Equation K n k} --- the free algebra over $K$ on the unknowns
    \texttt{inl l} and the constants \texttt{inr A}, $A\in M_n(K)$;
    \texttt{value X} --- evaluation, sending \texttt{inl l} to $X_l$ and
    \texttt{inr A} to $A$; \texttt{solutions S} --- the algebraic set $V(S)$.
  \item \texttt{EqNoetherian K n} --- the definition above;
    \texttt{ConjectureHolds} --- for every field $K$ (in \texttt{Type}) and
    every $n$, \texttt{EqNoetherian K n} $\leftrightarrow n\le2$.
  \item \texttt{generic}, \texttt{generic\_eval} --- the generic-matrix
    homomorphism and its compatibility with evaluation (by
    \texttt{FreeAlgebra.hom\_ext}).
  \item \texttt{eqNoetherian} --- Theorem~\ref{thm:main}, via Mathlib's
    Noetherianity of multivariate polynomial rings and
    \texttt{Submodule.fg\_span\_iff\_fg\_span\_finset\_subset}.
  \item \texttt{conjecture\_00000002058\_false} --- $\lnot\,$\texttt{ConjectureHolds},
    at $K=\mathbb Q$, $n=3$.
\end{itemize}

The toolchain is \texttt{leanprover/lean4:v4.33.1}, Mathlib is pinned to
\texttt{v4.33.1}, and \texttt{lake build} succeeds. The project contains no
\texttt{sorry}, no \texttt{admit} and no \texttt{native\_decide}, and
introduces no axioms:
\begin{center}
\texttt{\#print axioms conjecture\_00000002058\_false}
\end{center}
reports exactly \texttt{propext}, \texttt{Classical.choice} and
\texttt{Quot.sound}.

\section{Remarks}

The same argument shows that every finite-dimensional algebra over a field is
equationally Noetherian. Infinite ascending chains of ideals of
noncommutative polynomials (``commutator polynomials'') do exist in free
algebras. But equational Noetherianity concerns the solution sets of
equations, not ideals of the free algebra, and these stabilize because they are
algebraic sets in an affine space over $K$.

\begin{thebibliography}{9}
\bibitem{bmr} G.~Baumslag, A.~Myasnikov and V.~Remeslennikov, \emph{Algebraic
  geometry over groups I. Algebraic sets and ideal theory}, J. Algebra 219
  (1999), 16--79.
\end{thebibliography}

\end{document}
